\clearpage
\begin{trapbox}
\begin{center}\begin{tikzpicture}[font=\small,>=stealth]
\draw[gray!45,rounded corners=3pt] (0,0) rectangle (14.2,10);
\node at (7.1,9.35) {只保留平行还不够：把直线平移，线段比就改变了};

\node at (3.6,8.1) {过焦点：$d=c=5$};
\node at (10.6,8.1) {未过焦点：$d=4$};

\begin{scope}[shift={(0.5,5.9)},scale=.85]

\draw[teal!75!black,thick,samples=90,domain=-4.2:1.4,variable=\v] plot ({3*sqrt(1+\v*\v/16)},{\v});

\draw[gray] (1,0)--(6,0);
\draw[violet,dashed] (0,0)--(3.2,-64/15);
\draw[blue!75!black,very thick] (5,0)--(3.4,-2.1333333333333333);
\draw[orange!85!black,very thick] (3.4,-2.1333333333333333)--(2.5,-3.3333333333333335);
\fill (5,0) circle(2.8pt) node[above left=5pt] {$Q$};
\fill (3.4,-2.1333333333333333) circle(2.8pt) node[right=5pt] {$M$};
\fill (2.5,-3.3333333333333335) circle(2.8pt) node[below left=5pt] {$N$};
\fill (5,0) circle(2.8pt) node[above right=5pt] {$F$};
\end{scope}

\begin{scope}[shift={(7.5,5.9)},scale=.85]

\draw[teal!75!black,thick,samples=90,domain=-4.2:1.4,variable=\v] plot ({3*sqrt(1+\v*\v/16)},{\v});

\draw[gray] (1,0)--(6,0);
\draw[violet,dashed] (0,0)--(3.2,-64/15);
\draw[blue!75!black,very thick] (4,0)--(3.125,-1.1666666666666665);
\draw[orange!85!black,very thick] (3.125,-1.1666666666666665)--(2.0,-2.6666666666666665);
\fill (4,0) circle(2.8pt) node[above left=5pt] {$Q$};
\fill (3.125,-1.1666666666666665) circle(2.8pt) node[right=5pt] {$M$};
\fill (2.0,-2.6666666666666665) circle(2.8pt) node[below left=5pt] {$N$};
\fill (5,0) circle(2.8pt) node[above right=5pt] {$F$};
\end{scope}

\node at (3.6,1.9) {$Q=F,\quad QM/MN=16/9$};
\node at (10.6,1.9) {$Q\ne F,\quad QM/MN=7/9$};
\node at (7.1,.65) {两图同一双曲线、同一尺度、同一斜率；差别在于直线是否过焦点。};

\end{tikzpicture}
\end{center}
\textbf{先看对照：}只把直线平移，尽管斜率不变，截线比已经改变。
\begin{itemize}
\item \textbf{方向翻转：}原条件为
$\overrightarrow{FM}=\lambda\overrightarrow{MN}$（$\lambda>0$），
故 $\lambda=FM/MN$。改成 $\overrightarrow{MN}=\mu\overrightarrow{FM}$ 时，
$e=\sqrt{1+1/\mu}$，不是 $\sqrt{1+\mu}$。
\item \textbf{错把顺序写成 $F-N-M$：}本构型 $F-M-N$。代数上
$N_x=c/2$，$M_x=c/2+a^2/(2c)$，$F_x=c$，故 $c/2<x_M<c$。
\item \textbf{两条渐近线用成同一条：}过焦点的线平行一条渐近线，
$N$ 必须是与\textbf{另一条}渐近线的交点；两条平行线本身没有有限交点。
\item \textbf{没有过焦点就直接套：}本结论同时依赖“过焦点”和“平行渐近线”；
右图的平行线以 $Q=(4,0)$ 为起点，不经过 $F=(5,0)$，其 $QM/MN=7/9$，不等于 $16/9$。
\item \textbf{把 $e^2-1$ 误记为 $e-1$：}
$\lambda=b^2/a^2=(c^2-a^2)/a^2=\boxed{e^2-1}$。
\item \textbf{把 $\lambda$ 记成 $b/a$：}渐近线斜率绝对值是
$k=b/a$，而线段比为\textbf{平方} $\lambda=k^2$。
\item \textbf{误用椭圆的焦距关系：}这里
$\boxed{c^2=a^2+b^2}$，不是 $a^2-b^2$。
\item \textbf{只比较横坐标差却没说明共线：}因为三点在斜率固定的同一条直线上，
长度与横坐标差才有共同倍数 $c/a$，因此比值相同。
\item \textbf{把 $MN=a/2$ 误读成“$M$、$N$ 坐标与 $b$ 无关”：}
只有\textbf{长度}定值；两点坐标一般随 $b$ 变化。
\item \textbf{忘了原式是向量式：}$\lambda>0$ 表示两向量同向；
若题目出现负向量比例或不同点序，要先重画并按几何关系重新分析，不能原样套。
\item \textbf{角度记错：}$\tan\theta=b/a$ 中
$\theta$ 是渐近线对 $x$ 轴的锐角，不必等于两条渐近线所成的锐角。
\item \textbf{漏掉 $e>1$：}真正的非退化双曲线满足 $b>0$，
因此 $\lambda>0$、$e>1$；$e=1$ 仅是退化极限，不能作为本题实际取值。
\end{itemize}
\end{trapbox}
